% !TEX program = xelatex
% Статья для сборника НИТУ «МИСиС» (шаблон sadekovlar/misis-article-latex)

\documentclass{misisarticle}

% --- Классификационные коды ---
\grnti{34.39.19, 15.21.35, 28.23.25}
\asjc{2809, 2805, 1702}
\oecd{5010, 1020}

% --- Название ---
\title{Латерализация слуховой реакции: сравнительный анализ времени ответа при стимуляции левого и правого уха}

% --- Авторы ---
\misisauthor{Е.~С.~Горчакова}{кафедра инженерной кибернетики НИТУ «МИСиС»}{Москва, Россия}{kocinaliza367@gmail.com}

% --- Аннотация и ключевые слова ---
\abstract{%
  Проверяется, отличаются ли скорость и точность реакции на звук в зависимости от уха предъявления~--- левого или правого. Ранее аналогичная асимметрия была обнаружена в сигнале ЭЭГ: компонента P300 слухового вызванного потенциала различалась при стимуляции левого и правого уха. В настоящей работе исследуется её поведенческий аналог. Данные собираются самостоятельно: 10--12 участников выполняют слуховую oddball-задачу с монауральным предъявлением редких и частых тонов. Различия между ушами оцениваются парным непараметрическим тестом, а также бинарным классификатором, предсказывающим ухо предъявления по паттерну ответа; качество классификатора проверяется на участниках, не входивших в обучающую выборку.
}

\keywords{%
  латерализация слуха, время реакции, слуховая асимметрия, редкий звуковой сигнал, поведенческий эксперимент, машинное обучение
}

\begin{document}

\maketitle

% ============================================================
\section{Введение}

Латерализация~--- неравномерное распределение функций между полушариями головного мозга. Для слуховой системы она проявляется в том, что звук, поданный в левое и в правое ухо, обрабатывается неодинаково: слуховые пути в значительной мере перекрёстные, поэтому сигнал от каждого уха преимущественно поступает в противоположное полушарие~\cite{naatanen1987}, а полушария специализированы на разных характеристиках звука~\cite{kimura1961, kimura1964}.

В предшествующей научно-исследовательской работе автора~\cite{gorchakova2026} на данных слуховой oddball-парадигмы датасета MNE-sample~\cite{gramfort2013} был выполнен полный конвейер анализа вызванных потенциалов ЭЭГ. Пространственно-временной кластерный пермутационный тест~\cite{maris2007} выявил значимое различие ответов на левосторонний и правосторонний слуховой стимул в центро-теменной области на латентности 280--430~мс (p~<~0{,}05), что соответствует компоненте P300~\cite{polich2007}. Иными словами, асимметрия обработки звука была обнаружена на уровне сигнала мозга. Однако работа имела два ограничения: данные одного испытуемого и отсутствие поведенческих показателей.

Открытым остаётся вопрос, проявляется ли эта асимметрия в поведении человека~--- во времени и точности реакции на звук. Целью настоящей работы является проверка наличия поведенческой асимметрии при обнаружении редкого тона, поданного в левое или правое ухо. Для достижения цели решаются следующие задачи:
\begin{itemize}
  \item разработать и реализовать протокол поведенческого эксперимента с монауральной подачей звука;
  \item собрать собственный датасет на группе участников;
  \item сравнить время реакции и точность обнаружения между ушами статистическими методами;
  \item оценить, можно ли определить ухо предъявления по поведенческому ответу методами машинного обучения.
\end{itemize}

% ============================================================
\section{Анализ работ}

\subsection{Асимметрия слухового восприятия}

Классические исследования слуховой асимметрии выполнены в парадигме дихотического прослушивания, когда в оба уха одновременно подаются разные стимулы. Д.~Кимура показала преимущество правого уха при восприятии речевых стимулов~\cite{kimura1961} и, напротив, преимущество левого уха при восприятии мелодий~\cite{kimura1964}. Эти результаты легли в основу представления о специализации левого полушария на речи, а правого~--- на высотных и пространственных характеристиках звука. Современный обзор методологии дихотического прослушивания приведён в работе~\cite{westerhausen2019}; в ней подчёркивается, что эффект в классической постановке во многом обусловлен конкуренцией одновременно поданных стимулов.

\subsection{Латерализация восприятия простых звуков}

Для простых тональных стимулов картина менее определённа. В работе~\cite{sininger2012} при исследовании порогов различения частоты, интенсивности и пауз для тонов было обнаружено преимущество левого уха, тогда как ожидаемое преимущество правого уха для шумовых стимулов не подтвердилось; эффекты оказались слабыми. В работе~\cite{morita2023} показано, что асимметрия латерализации звука зависит от возраста слушателя и от формы нарастания сигнала. Таким образом, направление и величина асимметрии для простых звуков зависят от задачи и характеристик стимула, а поведенческие данные для монауральной подачи тонов в задаче обнаружения практически отсутствуют.

\subsection{Анализ биосигналов и поведения методами машинного обучения}

Компонента P300 широко используется в нейроинтерфейсах и исследованиях внимания~\cite{polich2007}. В студенческой работе~\cite{zhuravlev2022} сигналы ЭЭГ с волной P300 классифицировались нейронной сетью с точностью около 71\%. В работах сборника НИТУ «МИСиС» методы машинного обучения применяются к поведенческим и физиологическим проявлениям человека: предсказанию точки фиксации взгляда по видеопотоку~\cite{zhuravleva2026}, распознаванию эмоций по изображению лица~\cite{honer2024} и голосовому управлению по естественной речи~\cite{saveliev2025}. Эти работы показывают применимость моделей классификации к данным о поведении человека, однако задача декодирования уха предъявления по времени реакции в них не рассматривается.

Из анализа следует, что асимметрия хорошо изучена для речевых стимулов при одновременной подаче в оба уха, тогда как для простых тонов при монауральной подаче в задаче обнаружения данных мало и они противоречивы. Эта постановка и рассматривается в настоящей работе.

% ============================================================
\section{Постановка задачи}

Рассматривается слуховая oddball-задача: участник слышит последовательность частых стандартных тонов, среди которых изредка появляется целевой тон более высокой частоты. При обнаружении целевого тона участник нажимает клавишу. Каждый тон подаётся строго в одно ухо~--- левое или правое.

Для участника $s$ и уха $e \in \{L, R\}$ определяются медианное время реакции на целевые тоны $\mathrm{RT}_{s,e}$ и индекс чувствительности
\begin{equation}
  d'_{s,e} = z(H_{s,e}) - z(F_{s,e}),
  \label{eq:dprime}
\end{equation}
где $H_{s,e}$~--- доля обнаруженных целевых тонов, $F_{s,e}$~--- доля ложных тревог на стандартные тоны, $z(\cdot)$~--- функция, обратная функции стандартного нормального распределения.

Проверяется гипотеза $H_1$: разности $\Delta\mathrm{RT}_s = \mathrm{RT}_{s,R} - \mathrm{RT}_{s,L}$ и $\Delta d'_s = d'_{s,R} - d'_{s,L}$ в среднем по участникам отличны от нуля. Нулевая гипотеза $H_0$ предполагает отсутствие различий между ушами. Гипотеза двусторонняя, так как литература не даёт однозначного направления эффекта для простых тонов.

Дополнительно ставится задача бинарной классификации: по признакам отдельной пробы (тип стимула, наличие ответа, время реакции, номер пробы) предсказать ухо предъявления $e$. Качество классификатора на участниках, не входивших в обучающую выборку, значимо превышающее случайный уровень, будет свидетельствовать о том, что асимметрия устойчиво проявляется в поведении.

% ============================================================
\section{Описание методов}

\subsection{Стимулы и экспериментальная установка}

Эксперимент реализован на языке Python в среде PsychoPy~\cite{peirce2019}. Стимулы~--- синтезированные чистые тоны: стандартный частотой 1000~Гц и целевой частотой 1800~Гц, длительностью 150~мс, с линейным нарастанием и затуханием по 10~мс для устранения щелчков. Монауральная подача обеспечивается тем, что сигнал воспроизводится в одном канале стереонаушников, а второй канал заглушён. Для всех участников используются одни и те же проводные наушники и единая громкость, зафиксированная перед первой сессией.

\subsection{Протокол эксперимента}

В эксперименте участвуют 10--12 человек. Перед сессией участник даёт согласие на использование обезличенных данных, указывает ведущую руку и сообщает о жалобах на слух. Сессия включает тренировочный блок из 12 проб с обратной связью, который не входит в анализ, и основной блок. Параметры основного блока приведены в табл.~\ref{tab:protocol}. Порядок проб полностью рандомизирован по уху и типу стимула, что распределяет эффекты утомления и научения между условиями равномерно.

\begin{table}[htbp]
  \caption{Параметры основного блока}
  \label{tab:protocol}
  \centering
  \begin{tabular}{lc}
    \toprule
    \textbf{Параметр} & \textbf{Значение} \\
    \midrule
    Проб на участника & 500 \\
    Проб на каждое ухо & 250 \\
    Доля целевых проб & 20\% \\
    Межстимульный интервал & 1--2~с \\
    Длительность блока & $\approx$14~мин \\
    \bottomrule
  \end{tabular}
\end{table}

Для каждой пробы регистрируются ухо предъявления, тип стимула, факт ответа, время реакции от начала тона и правильность ответа. Ожидаемый объём датасета~--- 5000--6000 проб.

\subsection{Предобработка данных}

Ответ на целевой тон в интервале 100--1000~мс считается попаданием, его отсутствие~--- пропуском, ответ на стандартный тон~--- ложной тревогой. Ответы с временем реакции менее 100~мс считаются антиципационными и исключаются. Участники с долей попаданий ниже 70\% или долей ложных тревог выше 20\% исключаются из анализа. Для каждого участника и уха вычисляются медианное время реакции на попадания и индекс $d'$ по формуле~\eqref{eq:dprime}; медиана используется вместо среднего из-за асимметричности распределения времени реакции.

\subsection{Статистический анализ}

Разности $\Delta\mathrm{RT}_s$ и $\Delta d'_s$ сравниваются с нулём парным перестановочным тестом со сменой знака; в качестве дополнительной проверки применяется критерий Уилкоксона. Размер эффекта оценивается величиной Cohen's $d_z$, доверительный интервал разности~--- бутстрепом. На уровне отдельных проб строится линейная модель со смешанными эффектами для логарифма времени реакции с фиксированными эффектами уха и номера пробы и случайным эффектом участника. Поскольку выборка невелика, исследование рассматривается как пилотное: основное внимание уделяется направлению и величине эффекта.

\subsection{Методы машинного обучения}

Для задачи классификации уха предъявления используются логистическая регрессия и градиентный бустинг. Время реакции нормируется внутри каждого участника ($z$-преобразование), чтобы модель опиралась на различия между ушами, а не на индивидуальную скорость реакции; для проб без ответа время реакции заменяется пропуском с отдельным бинарным признаком. Валидация выполняется по схеме исключения одного участника (GroupKFold): модель всегда тестируется на человеке, данные которого не использовались при обучении, что исключает утечку индивидуальных особенностей. Качество оценивается сбалансированной точностью и ROC-AUC; значимость превышения случайного уровня проверяется перестановочным тестом с перемешиванием меток уха.

% ============================================================
\begin{thebibliography}{99}

\bibitem{naatanen1987} Näätänen R., Picton T. The N1 wave of the human electric and magnetic response to sound: a review and an analysis of the component structure. \textit{Psychophysiology}, vol.~24, no.~4, pp.~375--425, 1987.

\bibitem{kimura1961} Kimura D. Cerebral dominance and the perception of verbal stimuli. \textit{Canadian Journal of Psychology}, vol.~15, no.~3, pp.~166--171, 1961.

\bibitem{kimura1964} Kimura D. Left-right differences in the perception of melodies. \textit{Quarterly Journal of Experimental Psychology}, vol.~16, no.~4, pp.~355--358, 1964. DOI: 10.1080/17470216408416391.

\bibitem{gorchakova2026} Горчакова Е.~С. Анализ вызванных потенциалов ЭЭГ: методы предобработки, эпохирования и статистической оценки на примере слуховых ERP: отчёт о НИР. М.: НИТУ «МИСиС», 2026.

\bibitem{gramfort2013} Gramfort A., Luessi M., Larson E. et al. MEG and EEG data analysis with MNE-Python. \textit{Frontiers in Neuroscience}, vol.~7, art.~267, 2013. DOI: 10.3389/fnins.2013.00267.

\bibitem{maris2007} Maris E., Oostenveld R. Nonparametric statistical testing of EEG- and MEG-data. \textit{Journal of Neuroscience Methods}, vol.~164, no.~1, pp.~177--190, 2007. DOI: 10.1016/j.jneumeth.2007.03.024.

\bibitem{polich2007} Polich J. Updating P300: an integrative theory of P3a and P3b. \textit{Clinical Neurophysiology}, vol.~118, no.~10, pp.~2128--2148, 2007. DOI: 10.1016/j.clinph.2007.04.019.

\bibitem{westerhausen2019} Westerhausen R. A primer on dichotic listening as a paradigm for the assessment of hemispheric asymmetry. \textit{Laterality: Asymmetries of Body, Brain and Cognition}, vol.~24, no.~6, pp.~740--771, 2019. DOI: 10.1080/1357650X.2019.1598426.

\bibitem{sininger2012} Sininger Y.~S., Bhatara A. Laterality of basic auditory perception. \textit{Laterality: Asymmetries of Body, Brain and Cognition}, vol.~17, no.~2, pp.~129--149, 2012. DOI: 10.1080/1357650X.2010.541464.

\bibitem{morita2023} Morita K., Guo Y., Toi T. Age-related asymmetry in left--right ears of sound lateralization with respect to four different rise times. \textit{Frontiers in Neuroscience}, vol.~17, art.~1249119, 2023. DOI: 10.3389/fnins.2023.1249119.

\bibitem{zhuravlev2022} Журавлёв Д.~В., Резинченко А.~А. Обработка данных с волной P300 нейросетями // Информационные радиосистемы и радиотехнологии 2022: материалы науч.-техн. конф. Минск, 2022. С.~215--219.

\bibitem{zhuravleva2026} Журавлева А.~А. Регрессионная модель для предсказания точки фиксации взгляда пользователя на экране по данным RGB-видеопотока // Искусственный интеллект в промышленных, коммерческих, медицинских и финансовых приложениях: сб. ст. науч.-техн. семинара. Вып.~5. М.: НИТУ «МИСиС», 2026. С.~54--59.

\bibitem{honer2024} Хонер П.~Д. Классификация эмоций на лице человека при помощи компьютерного зрения // Искусственный интеллект в промышленных, коммерческих, медицинских и финансовых приложениях: сб. ст. науч.-техн. семинара. Вып.~2. М.: НИТУ «МИСиС», 2024. С.~140--146.

\bibitem{saveliev2025} Савельев Я.~С., Рябухин И.~А., Синельникова Т.~А. Применение больших языковых моделей в рамках голосового управления роботом-манипулятором посредством естественной речи // Искусственный интеллект в промышленных, коммерческих, медицинских и финансовых приложениях: сб. ст. науч.-техн. семинара. Вып.~3. М.: НИТУ «МИСиС», 2025. С.~116--119.

\bibitem{peirce2019} Peirce J., Gray J.~R., Simpson S. et al. PsychoPy2: experiments in behavior made easy. \textit{Behavior Research Methods}, vol.~51, no.~1, pp.~195--203, 2019. DOI: 10.3758/s13428-018-01193-y.

\end{thebibliography}

\end{document}
